% Nội dung chính. File này là VÍ DỤ cách dùng template (mục, mã nguồn, bảng, hình, công thức,
% trích dẫn) – xoá/thay bằng nội dung thật. Nội dung lớn có thể tách thành nhiều file trong
% sections/ rồi \input{sections/<tên>} trong main.tex.

\section{Ví dụ: đoạn mã cần kiểm thử}
\label{sec:code}

Đoạn mã dưới đây (ví dụ \texttt{daysInMonth} trong \citet{ammann2016introduction}) có nhiều câu
lệnh điều kiện lồng nhau. Mã nguồn đặt trong \texttt{lstlisting}, chọn ngôn ngữ bằng
\texttt{language=Java} (hoặc \texttt{Python}, \texttt{C}, \dots).

\begin{lstlisting}[language=Java, caption={Hàm \texttt{daysInMonth}}, label={lst:days}]
public static int daysInMonth(int m, int y) {
    if (m <= 0 || m > 12)            // p1: invalid month
        throw new IllegalArgumentException("Invalid month: " + m);
    if (m == 2) {                    // p2: February
        // p3: leap year
        if (y % 400 == 0 || (y % 4 == 0 && y % 100 != 0))
            return 29;
        else
            return 28;
    }
    if (m <= 7) {                    // p4
        if (m % 2 == 1)              // p5
            return 31;
        return 30;
    }
    if (m % 2 == 0)                  // p6
        return 31;
    return 30;
}
\end{lstlisting}

\section{Ví dụ: bảng test case}
\label{sec:tests}

Bảng dùng gói \texttt{booktabs} (\verb|\toprule|, \verb|\midrule|, \verb|\bottomrule|) và
không dùng đường kẻ dọc. Bảng dài hoặc bảng sinh tự động: đặt phần thân bảng trong thư mục
\texttt{tables/} rồi chèn bằng lệnh \verb|\inputtable{tables/<file>.tex}|.

\begin{table}[h]
\centering
\small
\caption{Một số test case cho \texttt{daysInMonth} (Mã~\ref{lst:days})}
\label{tab:tests}
\begin{tabular}{l l l l c}
\toprule
ID & Đầu vào $(m, y)$ & Kỳ vọng & Mục tiêu & Kết quả \\
\midrule
TC01 & $(0, 2024)$  & \texttt{IllegalArgumentException} & $p_1$ đúng (clause $c_1$) & \pass \\
TC02 & $(2, 2000)$  & 29 & $p_3$ đúng ($c_4$: chia hết cho 400) & \pass \\
TC03 & $(2, 1900)$  & 28 & $p_3$ sai ($c_6$ sai) & \pass \\
TC04 & $(8, 2023)$  & 31 & $p_4$ sai, $p_6$ đúng & \pass \\
\bottomrule
\end{tabular}
\end{table}

\section{Ví dụ: công thức, hình và trích dẫn}
\label{sec:misc}

Độ bao phủ nhánh được tính theo công thức
\begin{equation}
  \mathrm{BC} = \frac{\text{số nhánh đã thực thi}}{\text{tổng số nhánh}} \times 100\%.
  \label{eq:bc}
\end{equation}

Hình đặt trong \texttt{figures/} (PDF/PNG) và chèn như Hình~\ref{fig:logo}. Trích dẫn dạng số
\cite{fraser2011evosuite} hoặc dạng tác giả \citet{fraser2011evosuite}; khoá trích dẫn nằm trong
\texttt{refs.bib}. Phần còn thiếu đánh dấu bằng \verb|\todo{...}|, ví dụ: \todo{bổ sung số liệu
coverage đo bằng JaCoCo}.

\begin{figure}[h]
\centering
\includegraphics[height=2cm]{hust_logo.png}
\caption{Ví dụ chèn hình}
\label{fig:logo}
\end{figure}
